\documentclass[12pt,a4paper]{article}
\usepackage[utf8]{inputenc}
\usepackage[T1]{fontenc}
\usepackage[french]{babel}
\usepackage{geometry}
\geometry{margin=2.5cm}

\title{L'Histoire de l'Intelligence Artificielle}
\author{Lily AUBERT, Maelys EL ARAARI, Baptiste MASSON}
\date{Mars 2026}

\begin{document}

\maketitle

\section{Les origines et les prémisses}

L'IA ne sort pas du néant ; elle s'appuie sur des travaux séculaires et des mythes anciens. 

\begin{itemize}
    \item \textbf{Mythes et illusions :} En 1769, le \textbf{Turc mécanique} de Von Kempelen simulait déjà un automate capable de jouer aux échecs, bien qu'il s'agissait d'une supercherie.
    \item \textbf{Bases théoriques :} Elle se nourrit de la logique aristotélicienne, de la machine à calculer de Pascal, de la cybernétique de Wiener et de la théorie des automates de Turing.
    \item \textbf{Le Test de Turing :} Alan Turing a posé les bases de la réflexion sur la capacité des machines à simuler l'intelligence humaine.
\end{itemize}

\section{1956 : La naissance officielle}

Le terme \textbf{« Intelligence Artificielle »} a été introduit par \textbf{John McCarthy} lors du séminaire du \textbf{Dartmouth College} en août 1956. Les « pères fondateurs » présents incluaient également Marvin Minsky, Claude Shannon et Nathaniel Rochester.

\section{Les cycles de l'IA : Étés et Hivers}

L'histoire de l'IA est marquée par des périodes d'enthousiasme (étés) suivies de phases de désillusion et de baisse de financement (hivers).

\begin{itemize}
    \item \textbf{Premier âge d'or (1956-1974) :} Création de programmes comme \textit{Logic Theorist} (1955), \textit{ELIZA} (1966) — le premier robot conversationnel — et \textit{SHRDLU}.
    \item \textbf{Premier hiver (1974-1980) :} Déclenché notamment par le \textbf{rapport Lighthill} en 1973, qui critiquait le manque de résultats concrets.
    \item \textbf{Regain et systèmes experts (1980-1987) :} Succès de systèmes comme \textit{XCON} et projets comme la 5ème génération japonaise.
    \item \textbf{Second hiver (1987-1993) :} Nouvelle période de stagnation.
    \item \textbf{L'essor moderne (Depuis 2002) :} Porté par l'augmentation de la puissance de calcul (supercalculateurs comme J. Zay) et l'accès massif aux données (Web 2.0).
\end{itemize}

\section{Grands jalons et succès médiatiques}

L'évolution récente est marquée par des victoires symboliques de la machine sur l'humain dans des jeux complexes :

\begin{itemize}
    \item \textbf{1997 :} \textbf{Deep Blue} bat Garry Kasparov aux échecs.
    \item \textbf{2011 :} \textbf{Watson} (IBM) gagne au jeu \textit{Jeopardy!} en s'attaquant au langage naturel.
    \item \textbf{2016 :} \textbf{AlphaGo} bat Lee Sedol, marquant le succès des réseaux de neurones et de l'apprentissage profond (\textit{deep learning}).
\end{itemize}

\section{L'ère contemporaine}

Depuis 2017, l'IA a franchi de nouvelles étapes avec l'apparition des \textbf{Transformers} (2017), de \textbf{BERT} (2018), puis d'outils génératifs comme \textbf{DALL-E} (2021) et \textbf{ChatGPT} (2022). Aujourd'hui, l'IA est au cœur de stratégies nationales massives, à l'image du Plan Villani en France ou du projet Stargate aux États-Unis.

\end{document}
